\chapter{Introduction}

\begin{aside}
\textbf{Sources.} The module list and ordering in these notes follow the archived CS 350 term
pages on \code{student.cs.uwaterloo.ca/\textasciitilde cs350/}, specifically the course outline and
reading-materials pages of \textbf{Winter 2018} (modules: Introduction, Threads, Synchronization,
Processes/Kernel/System Calls, Virtual Memory, Scheduling, Devices and I/O, File Systems; the
``Interprocess Communication and Networking'' module was listed as \emph{not covered}) and
\textbf{Fall 2020} (same modules, with a new \emph{Virtual Machines} module at the end). The
recommended (optional) textbook in both terms was Arpaci-Dusseau, \emph{Operating Systems: Three
Easy Pieces} (OSTEP); older terms used Silberschatz, Galvin and Gagne. Assignments were done in
OS/161 on the System/161 MIPS simulator.
\end{aside}

\section{What is an operating system?}

There is no single crisp definition; it helps to look at the OS from three viewpoints.

\begin{description}
  \item[Application view] The OS provides an \emph{execution environment} for running
    programs. It gives each program the resources it needs (CPU time, memory, storage) and an
    \emph{interface} to those resources that is more convenient than the bare hardware:
    files instead of disk sectors, virtual address spaces instead of physical RAM,
    threads instead of a single CPU.
  \item[System view] The OS is a \emph{resource manager}. It allocates resources among many
    concurrently running programs, isolates them from each other (one buggy program should
    not crash another), and enforces policy (who may access which file, who runs next).
  \item[Implementation view] The OS is a \emph{concurrent, real-time program}. Concurrency
    is inherent: many programs, devices and interrupts are all active at once. Real-time
    because hardware events (a network packet, a disk completion) must be handled promptly.
\end{description}

\begin{definition}[Kernel]
The \textbf{kernel} is the part of the OS that runs in privileged (kernel / supervisor) mode
and responds to system calls, interrupts and exceptions. The \textbf{operating system} as a whole
also includes user-level parts: system libraries, daemons/servers, shells, utilities.
\end{definition}

\begin{important}[Privilege]
Hardware supports at least two execution modes. In \emph{kernel mode} the CPU may execute
privileged instructions (halt, change the page table base / TLB, disable interrupts, talk to
devices). In \emph{user mode} such instructions trap. Applications run in user mode, so the only
way for an application to use hardware is to ask the kernel through a \emph{system call}.
\end{important}

\section{Layers of a computer system}

\begin{center}
\begin{tikzpicture}[box/.style={draw,text width=7.6cm,minimum height=0.8cm,align=center,font=\small}]
  \node[box] (u) {User programs (shell, editors, compilers, \ldots)};
  \node[box,below=0pt of u] (l) {Libraries (libc: \code{printf}, \code{malloc}, wrappers)};
  \node[box,below=0pt of l,fill=blue!8] (k) {Kernel: processes, VM, scheduler, FS, drivers};
  \node[box,below=0pt of k] (h) {Hardware: CPU, MMU/TLB, memory, devices};
  \draw[-{Stealth},thick] ($(l.north east)+(0.3,0)$) -- node[right,font=\small]{system calls} ($(k.south east)+(0.3,0.1)$);
  \draw[-{Stealth},thick] ($(h.south west)+(-0.3,0)$) -- node[left,font=\small]{interrupts, exceptions} ($(k.north west)+(-0.3,-0.1)$);
\end{tikzpicture}
\end{center}

The kernel sits between the user programs and the hardware. Control enters the kernel in exactly
three ways:
\begin{enumerate}
  \item \textbf{System calls} --- deliberate requests from a user program (via a trap
    instruction such as \code{syscall} on MIPS).
  \item \textbf{Exceptions} --- errors or conditions caused by the running instruction
    (divide by zero, illegal address, TLB miss).
  \item \textbf{Interrupts} --- asynchronous signals from devices (timer tick, disk done).
\end{enumerate}
After handling, the kernel either returns to the interrupted program or switches to another.

\section{OS abstractions}

A large part of the course is about how the OS implements a few key abstractions.

\begin{center}
\begin{tabular}{@{}lll@{}}
\toprule
Abstraction & Hides / virtualizes & Module \\
\midrule
Thread & the CPU (sequential execution) & Threads, Scheduling \\
Process & address space + threads + resources & Processes \\
Address space & physical memory & Virtual memory \\
File / directory & disk blocks & File systems \\
Device driver interface & device registers & Devices and I/O \\
Virtual machine & the entire hardware & Virtual machines \\
\bottomrule
\end{tabular}
\end{center}

\begin{aside}
OSTEP organizes the same material as three ``easy pieces'': \emph{virtualization} (CPU and
memory), \emph{concurrency} (threads, locks, condition variables), and \emph{persistence}
(devices and file systems). CS 350 covers all three, with concurrency first because writing a
kernel requires synchronization from day one.
\end{aside}

\section{Monolithic kernels vs microkernels}

\begin{description}
  \item[Monolithic kernel] Everything (file systems, drivers, networking, VM, scheduling) runs
    in kernel mode in one address space. Fast (procedure calls between subsystems) but a bug
    anywhere can crash the system. Linux, traditional Unix, OS/161.
  \item[Microkernel] Only a minimum runs in kernel mode (address spaces, threads, IPC);
    file systems, drivers, etc.\ run as user-level server processes. More robust and modular,
    but more IPC overhead. Mach, L4, QNX, MINIX 3.
  \item[Hybrid] Mixtures, e.g.\ Windows NT and macOS (XNU = Mach + BSD).
\end{description}

\section{OS/161 and System/161}

The assignments use \textbf{OS/161}, a small teaching operating system originally from Harvard,
which runs on \textbf{System/161}, a simulator of a MIPS R3000 based machine (\code{sys161}).

\begin{itemize}
  \item The kernel is written in C (with a little MIPS assembly for low-level entry/exit and
    context switching).
  \item Source tree highlights: \code{kern/thread} (threads, scheduler, synchronization
    primitives in \code{synch.c}), \code{kern/arch/mips} (trap handling, \code{switchframe},
    TLB code), \code{kern/syscall}, \code{kern/vm} (initially the very simple \code{dumbvm}),
    \code{kern/vfs} and \code{kern/fs} (SFS), plus user-level programs and \code{libc} under
    \code{user/}.
  \item Debugging is done by attaching \code{cs350-gdb} to the simulator (\code{sys161 -w}).
\end{itemize}

Typical assignment progression (2016--2021):
\begin{enumerate}[label=\textbf{A\arabic*}]
  \setcounter{enumi}{-1}
  \item Getting familiar with the code base, building and running the kernel, \code{gdb}.
  \item Implementing \textbf{locks} and \textbf{condition variables}, then solving a
    synchronization problem (e.g.\ a traffic intersection or cat/mouse simulation).
  \item[\textbf{A2a}] Process system calls: \code{fork}, \code{getpid}, \code{waitpid}, \code{\_exit}.
  \item[\textbf{A2b}] \code{execv} with argument passing.
  \item[\textbf{A3}] Virtual memory: TLB replacement, read-only text segment, replacing
    \code{dumbvm} with a page-granular physical memory allocator (coremap).
\end{enumerate}

\begin{lstlisting}[style=bash]
# build and run OS/161 (typical workflow)
cd os161-1.99/kern/compile/ASST2
bmake depend && bmake && bmake install
cd ~/cs350-os161/root
sys161 kernel-ASST2          # boots to the OS/161 menu
sys161 -w kernel-ASST2       # wait for debugger
cs350-gdb kernel             # in another terminal: target remote unix:.sockets/gdb
\end{lstlisting}

\section{Course roadmap}

\begin{enumerate}
  \item Threads and concurrency --- how one CPU appears to run many threads.
  \item Synchronization --- making concurrent threads cooperate safely.
  \item Processes and the kernel --- system calls, exceptions, interrupts, user/kernel boundary.
  \item Virtual memory --- address translation, paging, TLBs, page replacement.
  \item Scheduling --- choosing which thread runs next.
  \item Devices and I/O --- talking to hardware, disks and SSDs.
  \item File systems --- persistent storage, i-nodes, crash consistency.
  \item Virtual machines --- running whole operating systems on top of a hypervisor.
\end{enumerate}

\section*{Summary}
\begin{itemize}
  \item An OS provides an execution environment (abstractions + interface), manages and
    isolates resources, and is itself a concurrent real-time program.
  \item The kernel runs in privileged mode; it is entered only through system calls,
    exceptions and interrupts.
  \item Key abstractions: threads, processes, address spaces, files, virtual machines.
  \item CS 350 assignments build process management and virtual memory into OS/161, a small
    monolithic kernel running on a simulated MIPS R3000.
\end{itemize}
